% =========================================================
%  ppt/03-gerak-vector2-delta-time.tex
%  Materi 3 — Gerak Objek, Vector2, Input, dan Delta Time
%  Bab V buku ajar
% =========================================================
\documentclass[aspectratio=169,10pt]{beamer}
% =========================================================
%  ppt/config/packages.tex
%  Paket bersama untuk seluruh presentasi kuliah
%  Pengantar Pemrograman Game
% =========================================================

\usepackage[T1]{fontenc}
\usepackage[utf8]{inputenc}
\usepackage[indonesian]{babel}
\usepackage{lmodern}
\usepackage{textcomp}
\usepackage{amsmath}
\usepackage{amssymb}
\usepackage{graphicx}
\usepackage{booktabs}
\usepackage{tabularx}
\usepackage{array}
\usepackage{multicol}
\usepackage{listings}
\usepackage{upquote}
\usepackage{adjustbox}
\usepackage{tcolorbox}
\usepackage{tikz}
\usetikzlibrary{
  arrows.meta,
  positioning,
  shapes.geometric,
  shapes.misc,
  calc,
  fit,
  backgrounds,
  matrix,
  chains,
  decorations.pathreplacing,
  patterns
}

\tcbuselibrary{skins}

% Jangan memuat geometry: ukuran halaman diatur oleh beamer.

% =========================================================
%  ppt/config/theme.tex
%  Satu tema Beamer untuk keempat belas presentasi
%  Pengantar Pemrograman Game
% =========================================================

% --------------------------------------------------------- warna
\definecolor{ubNavy}{HTML}{1F3864}
\definecolor{ubAccent}{HTML}{2E75B6}
\definecolor{ubHead}{HTML}{D9E2F3}
\definecolor{ubSoft}{HTML}{F2F6FC}
\definecolor{ubLine}{HTML}{BFBFBF}
\definecolor{ubText}{HTML}{1F2937}
\definecolor{ubMuted}{HTML}{5A6678}
\definecolor{ubGreen}{HTML}{1E7B45}
\definecolor{ubGreenBg}{HTML}{F1F8F2}
\definecolor{ubRed}{HTML}{B3261E}
\definecolor{ubRedBg}{HTML}{FDF2F1}
\definecolor{ubWarn}{HTML}{A35A05}
\definecolor{ubWarnBg}{HTML}{FFF8EC}
\definecolor{ubGray}{HTML}{F6F6F6}
\definecolor{ubCodeBg}{HTML}{10233F}
\definecolor{ubCodeText}{HTML}{E8EEF7}
\definecolor{ubCodeKey}{HTML}{8FBCF0}
\definecolor{ubCodeCom}{HTML}{8FB37A}
\definecolor{ubCodeStr}{HTML}{D9A46B}
\definecolor{ubCodeNum}{HTML}{7E99BE}

% Nama warna yang dipakai gambar TikZ buku ajar.
\definecolor{ubPrimary}{HTML}{1F3864}
\definecolor{ubAlt}{HTML}{F2F6FC}

% --------------------------------------------------------- dasar
\usetheme{default}
\usefonttheme{professionalfonts}
\setbeamertemplate{navigation symbols}{}
\setbeamerfont{frametitle}{size=\Large,series=\bfseries}
\setbeamerfont{framesubtitle}{size=\small}
\setbeamerfont{footline}{size=\tiny}
\setbeamerfont{block title}{size=\small,series=\bfseries}
\setbeamerfont{block body}{size=\small}
\setbeamerfont{itemize/enumerate body}{size=\small}
\setbeamerfont{itemize/enumerate subbody}{size=\footnotesize}

\setbeamersize{text margin left=0.60cm,text margin right=0.60cm}
\uchyph=0

\setbeamercolor{normal text}{fg=ubText,bg=white}
\setbeamercolor{structure}{fg=ubAccent}
\setbeamercolor{frametitle}{fg=ubNavy,bg=white}
\setbeamercolor{title}{fg=ubNavy}
\setbeamercolor{footline}{fg=ubMuted,bg=white}
\setbeamercolor{itemize item}{fg=ubAccent}
\setbeamercolor{itemize subitem}{fg=ubAccent}
\setbeamercolor{block title}{fg=white,bg=ubNavy}
\setbeamercolor{block body}{fg=ubText,bg=ubSoft}
\setbeamercolor{block title alerted}{fg=white,bg=ubWarn}
\setbeamercolor{block body alerted}{fg=ubText,bg=ubWarnBg}
\setbeamercolor{block title example}{fg=white,bg=ubGreen}
\setbeamercolor{block body example}{fg=ubText,bg=ubGreenBg}
\setbeamercolor{caption name}{fg=ubMuted}

\setbeamertemplate{blocks}[rounded][shadow=false]
\setbeamertemplate{itemize item}{\raisebox{0.10em}{\color{ubAccent}\scriptsize$\blacksquare$}}
\setbeamertemplate{itemize subitem}{\raisebox{0.06em}{\color{ubAccent}\tiny$\blacktriangleright$}}
\setbeamertemplate{enumerate item}{%
  \makebox[1.3em][l]{\color{ubAccent}\insertenumlabel.}}

\addtobeamertemplate{itemize/enumerate body begin}{%
  \setlength{\itemsep}{0.42em}\setlength{\parskip}{0pt}}{}

% --------------------------------------------------------- judul frame
\setbeamertemplate{frametitle}[default][left]

% --------------------------------------------------------- kaki
\newcommand{\materino}{00}
\newcommand{\materifoot}{Pengantar Pemrograman Game \textperiodcentered\ Materi \materino}

\setbeamertemplate{footline}{%
  \begin{beamercolorbox}[wd=\paperwidth,ht=0.015cm,dp=0cm]{footline}%
    \color{ubHead}\rule{\paperwidth}{0.4pt}%
  \end{beamercolorbox}%
  \begin{beamercolorbox}[wd=\paperwidth,ht=0.40cm,dp=0.14cm,%
                         leftskip=0.60cm,rightskip=0.60cm]{footline}%
    \usebeamerfont{footline}\color{ubMuted}\materifoot\hfill\insertframenumber%
  \end{beamercolorbox}%
}

% --------------------------------------------------------- makro buku ajar
\providecommand{\code}[1]{\texttt{\detokenize{#1}}}
\providecommand{\istilah}[1]{\emph{#1}}
\providecommand{\CPMK}[1]{CPMK\,#1}
\providecommand{\CPL}[1]{CPL\,#1}
\providecommand{\SUBCPMK}[1]{Sub-CPMK\,#1}

% --------------------------------------------------------- gaya kode
\lstdefinelanguage{CSharp}{
  morekeywords={abstract,as,base,bool,break,byte,case,catch,char,class,const,
    continue,decimal,default,delegate,do,double,else,enum,event,explicit,
    extern,false,finally,fixed,float,for,foreach,goto,if,implicit,in,int,
    interface,internal,is,lock,long,namespace,new,null,object,operator,out,
    override,params,private,protected,public,readonly,ref,return,sbyte,
    sealed,short,sizeof,stackalloc,static,string,struct,switch,this,throw,
    true,try,typeof,uint,ulong,unchecked,unsafe,ushort,using,virtual,void,
    volatile,while,var,get,set},
  sensitive=true,
  morecomment=[l]{//},
  morecomment=[s]{/*}{*/},
  morestring=[b]",
}

\lstdefinestyle{kode}{
  basicstyle=\ttfamily\footnotesize\color{ubCodeText},
  keywordstyle=\color{ubCodeKey}\bfseries,
  commentstyle=\color{ubCodeCom}\itshape,
  stringstyle=\color{ubCodeStr},
  numberstyle=\tiny\color{ubCodeNum},
  numbers=left,
  numbersep=7pt,
  backgroundcolor=\color{ubCodeBg},
  frame=none,
  showstringspaces=false,
  breaklines=false,
  tabsize=4,
  columns=fullflexible,
  keepspaces=true,
  xleftmargin=12pt,
  xrightmargin=2pt,
  aboveskip=0pt,
  belowskip=0pt,
  upquote=true,
}
\lstset{style=kode}

\newcommand{\kodejelas}[1]{%
  \vspace{0.16cm}%
  {\footnotesize\color{ubMuted}#1\par}%
}
\newcommand{\sumber}[1]{%
  \vspace{0.08cm}%
  {\scriptsize\color{ubMuted}\itshape #1\par}%
}
\newcommand{\awalan}[1]{%
  {\small\color{ubMuted}#1\par}\vspace{0.18cm}%
}

% --------------------------------------------------------- kotak
\newtcolorbox{kotak}[1][]{enhanced, colback=ubSoft, colframe=ubAccent,
  boxrule=0pt, leftrule=2.2pt, arc=1pt, left=6pt, right=6pt, top=4pt,
  bottom=4pt, fontupper=\footnotesize, #1}
\newtcolorbox{peringatan}[1][]{enhanced, colback=ubWarnBg, colframe=ubWarn,
  boxrule=0pt, leftrule=2.2pt, arc=1pt, left=6pt, right=6pt, top=4pt,
  bottom=4pt, fontupper=\footnotesize, #1}
\newtcolorbox{kiat}[1][]{enhanced, colback=ubGreenBg, colframe=ubGreen,
  boxrule=0pt, leftrule=2.2pt, arc=1pt, left=6pt, right=6pt, top=4pt,
  bottom=4pt, fontupper=\footnotesize, #1}

% --------------------------------------------------------- sampul
% \deckcover{nomor}{minggu}{judul}{sub-cpmk}{teknologi}{sumber}{blok}
\newcommand{\deckcover}[7]{%
\begin{frame}[plain]
\begin{tikzpicture}[remember picture, overlay]
  \fill[ubNavy] (current page.south west) rectangle ([xshift=6.00cm]current page.north west);
  \fill[ubAccent] ([xshift=6.00cm]current page.south west) rectangle ([xshift=6.08cm]current page.north west);

  \node[anchor=north west, text width=4.72cm, align=left, inner sep=0pt]
    at ($(current page.north west)+(0.66cm,-0.60cm)$) {%
      {\color{white}\bfseries\fontsize{10.2}{12.2}\selectfont
       Universitas Bale Bandung}\\[4pt]
      {\color{ubHead}\fontsize{7.6}{9.4}\selectfont
       Fakultas Teknologi Informasi}\\[1pt]
      {\color{ubHead}\fontsize{7.6}{9.4}\selectfont
       Program Studi Teknik Informatika}\\[7pt]
      {\color{ubAccent}\rule{4.0cm}{0.5pt}}%
    };

  \node[anchor=north west, text width=4.72cm, align=left, inner sep=0pt]
    at ($(current page.north west)+(0.66cm,-2.72cm)$) {%
      {\color{ubAccent}\fontsize{6.4}{7.8}\selectfont MATA KULIAH}\\[3pt]
      {\color{white}\bfseries\fontsize{9.8}{12.0}\selectfont
       Pengantar Pemrograman Game}\\[3pt]
      {\color{ubHead}\fontsize{7.6}{9.4}\selectfont
       MWCA009 \textperiodcentered\ 3 SKS}\\[8pt]
      {\color{ubHead}\fontsize{7.6}{9.4}\selectfont
       Semester Ganjil 2026/2027}%
    };

  \node[anchor=south west, text width=4.72cm, align=left, inner sep=0pt]
    at ($(current page.south west)+(0.66cm,0.52cm)$) {%
      {\color{ubAccent}\fontsize{6.4}{7.8}\selectfont DOSEN PENGAMPU}\\[3pt]
      {\color{white}\fontsize{8.4}{10.4}\selectfont
       Cecep Suwanda, S.Si., M.Kom.}%
    };

  \node[anchor=north west, text width=8.95cm, align=left, inner sep=0pt]
    at ($(current page.north east)+(-9.40cm,-0.92cm)$) {%
      {\color{ubAccent}\bfseries\fontsize{8.2}{10.2}\selectfont
       MATERI #1\ \ \textperiodcentered\ \ MINGGU #2}\\[8pt]
      {\color{ubNavy}\bfseries\fontsize{19.4}{23.4}\selectfont
       \hyphenpenalty=10000 \exhyphenpenalty=10000 \raggedright #3}%
    };

  \node[anchor=north west, text width=8.95cm, align=left, inner sep=0pt]
    at ($(current page.north east)+(-9.40cm,-4.30cm)$) {%
      \renewcommand{\arraystretch}{1.30}%
      \begin{tabular}{@{}>{\color{ubMuted}\fontsize{7.3}{9.0}\selectfont}p{1.85cm}
                        >{\color{ubText}\fontsize{8.1}{10.0}\selectfont}p{6.7cm}@{}}
        Sub-CPMK      & #4 \\
        Teknologi     & #5 \\
        Sumber materi & #6 \\
        Blok          & #7 \\
      \end{tabular}%
    };
\end{tikzpicture}
\end{frame}%
}

% --------------------------------------------------------- pembatas bagian
\newcommand{\sectionframe}[3]{%
\begin{frame}[plain]
\begin{tikzpicture}[remember picture, overlay]
  \fill[ubNavy] (current page.south west) rectangle (current page.north east);
  \fill[ubAccent] (current page.north west) rectangle ([xshift=0.10cm]current page.south west);
  \node[anchor=west, inner sep=0pt]
    at ($(current page.west)+(1.00cm,1.45cm)$)
    {\color{ubNavy!52!white}\bfseries\fontsize{58}{58}\selectfont #1};
  \node[anchor=west, inner sep=0pt, text width=11.6cm, align=left]
    at ($(current page.west)+(1.05cm,-0.55cm)$) {%
      {\color{ubAccent}\bfseries\fontsize{8.0}{10.0}\selectfont BAGIAN}\\[3pt]
      {\color{white}\bfseries\fontsize{20.5}{24.5}\selectfont #2}\\[5pt]
      {\color{ubHead}\small #3}%
    };
  \node[anchor=south east, inner sep=0pt]
    at ($(current page.south east)+(-0.60cm,0.32cm)$)
    {\color{ubHead}\fontsize{7.4}{9}\selectfont #1 / 6
     \ \ \textperiodcentered\ \ \insertframenumber};
\end{tikzpicture}
\end{frame}%
}

% --------------------------------------------------------- peta pertemuan
\newcommand{\blokpeta}[6]{%
  \ifnum#2=#3
    \node[rounded corners=3pt, fill=ubNavy, draw=ubNavy, line width=0.5pt,
          text width=3.02cm, minimum height=1.62cm, align=left, inner sep=5pt]
      at (#1,0) {%
        {\color{ubAccent}\bfseries\fontsize{6.8}{8}\selectfont #4}\\[1pt]
        {\color{white}\bfseries\fontsize{7.6}{9}\selectfont #5}\\[2pt]
        {\color{ubHead}\fontsize{6.4}{7.6}\selectfont #6}};
  \else
    \node[rounded corners=3pt, fill=ubSoft, draw=ubHead, line width=0.5pt,
          text width=3.02cm, minimum height=1.62cm, align=left, inner sep=5pt]
      at (#1,0) {%
        {\color{ubAccent}\bfseries\fontsize{6.8}{8}\selectfont #4}\\[1pt]
        {\color{ubNavy}\bfseries\fontsize{7.6}{9}\selectfont #5}\\[2pt]
        {\color{ubMuted}\fontsize{6.4}{7.6}\selectfont #6}};
  \fi
}

% \petapertemuan{blok aktif}{sebelumnya}{sesudahnya}{keluaran}
% \petapertemuan{blok aktif}{sebelumnya}{pertemuan ini}{sesudahnya}{keluaran}
\newcommand{\petapertemuan}[5]{%
\begin{frame}{Peta Pertemuan}
{\scriptsize\color{ubMuted}Kedudukan materi dalam empat blok perkuliahan}\par
\vspace{0.10cm}
\centering
\begin{tikzpicture}
  \blokpeta{-5.58}{1}{#1}{BLOK I}{Materi 1--5}{Game loop, vektor,\\fisika, collision}
  \blokpeta{-1.86}{2}{#1}{BLOK II}{Materi 6--9}{AI NPC: chase, FSM,\\pathfinding A*}
  \blokpeta{1.86}{3}{#1}{BLOK III}{Materi 10--11}{Unity: scene,\\level, physics}
  \blokpeta{5.58}{4}{#1}{BLOK IV}{Materi 12--14}{Aset visual, audio,\\UI, playtest}
\end{tikzpicture}
\par\vspace{0.30cm}
\begin{columns}[T,onlytextwidth]
  \begin{column}{0.315\textwidth}
    \begin{beamercolorbox}[wd=\linewidth,sep=4pt,rounded=true]{block body}
      {\color{ubMuted}\fontsize{6.6}{8}\selectfont SEBELUMNYA}\\[2pt]
      {\footnotesize #2}
    \end{beamercolorbox}
  \end{column}
  \begin{column}{0.315\textwidth}
    \begin{beamercolorbox}[wd=\linewidth,sep=4pt,rounded=true]{block title}
      {\color{ubHead}\fontsize{6.6}{8}\selectfont PERTEMUAN INI}\\[2pt]
      {\footnotesize #3}
    \end{beamercolorbox}
  \end{column}
  \begin{column}{0.315\textwidth}
    \begin{beamercolorbox}[wd=\linewidth,sep=4pt,rounded=true]{block body}
      {\color{ubMuted}\fontsize{6.6}{8}\selectfont SESUDAHNYA}\\[2pt]
      {\footnotesize #4}
    \end{beamercolorbox}
  \end{column}
\end{columns}
\par\vspace{0.26cm}
{\footnotesize\color{ubText}\textbf{Keluaran pertemuan.} #5\par}
\end{frame}%
}

% --------------------------------------------------------- penutup
\newcommand{\penutup}[2]{%
\begin{frame}[plain]
\begin{tikzpicture}[remember picture, overlay]
  \fill[ubNavy] (current page.south west) rectangle (current page.north east);
  \fill[ubAccent] (current page.north west) rectangle ([xshift=0.10cm]current page.south west);

  \node[anchor=north west, text width=13.6cm, align=left, inner sep=0pt]
    at ($(current page.north west)+(1.00cm,-0.95cm)$) {%
      {\color{ubAccent}\bfseries\fontsize{8.2}{10}\selectfont ASESMEN MINGGU INI}\\[6pt]
      {\color{white}\bfseries\fontsize{24}{28}\selectfont #1}%
    };

  \node[anchor=north west, text width=13.6cm, align=left, inner sep=0pt]
    at ($(current page.north west)+(1.00cm,-4.55cm)$) {%
      {\color{ubNavy!0!ubHead}\rule{13.4cm}{0.4pt}}\\[7pt]
      {\color{ubAccent}\bfseries\fontsize{7.4}{9}\selectfont
       PENGAIT KE MATERI BERIKUTNYA}\\[4pt]
      {\color{ubHead}\fontsize{10.6}{13.2}\selectfont #2}%
    };

  \node[anchor=south east, inner sep=0pt]
    at ($(current page.south east)+(-0.60cm,0.32cm)$)
    {\color{ubHead}\fontsize{7.4}{9}\selectfont \insertframenumber};
\end{tikzpicture}
\end{frame}%
}



\renewcommand{\materino}{03}

\begin{document}

\deckcover{03}{3}
  {Gerak Objek, Vector2, Input, dan Delta Time}
  {1.2}
  {Python / Pygame}
  {Bab V buku ajar}
  {Blok I \textperiodcentered\ Materi 1--5}

% --------------------------------------------------------- capaian
\begin{frame}{Capaian Pertemuan}
\setbeamerfont{itemize/enumerate body}{size=\scriptsize}
\awalan{Sub-CPMK dan indikator dari RPS, serta tujuan Bab V buku ajar.}
\begin{columns}[T,onlytextwidth]
  \begin{column}{0.52\textwidth}
    \begin{beamercolorbox}[wd=\linewidth,sep=5pt,rounded=true]{block title}
      {\fontsize{6.8}{8}\selectfont\color{ubHead}SUB-CPMK 1.2}\\[2pt]
      {\footnotesize Merancang dan mengimplementasikan gerak objek berbasis
       \emph{delta time} dan vektor pada Pygame.}
    \end{beamercolorbox}
    \vspace{0.16cm}

    {\fontsize{7}{8.5}\selectfont\bfseries\color{ubAccent}INDIKATOR}\par
    \vspace{0.06cm}
    \begin{itemize}
      \item Menulis kode gerak menggunakan \texttt{pygame.Vector2} dan
            \emph{delta time}.
      \item Menunjukkan kecepatan objek tetap konsisten pada 30 \emph{FPS} dan
            60 FPS.
      \item Menyinkronkan posisi bertipe \emph{float} dengan \emph{Rect}
            (\texttt{pygame.Rect}) bertipe bilangan bulat.
    \end{itemize}
  \end{column}
  \begin{column}{0.45\textwidth}
    {\fontsize{7}{8.5}\selectfont\bfseries\color{ubAccent}TUJUAN BAB V}\par
    \vspace{0.08cm}
    \begin{itemize}
      \item Menyatakan posisi dan kecepatan sebagai vektor dengan
            \texttt{pygame.Vector2}.
      \item Membaca masukan yang berlangsung selama tombol ditahan.
      \item Menjelaskan mengapa vektor arah perlu dinormalkan.
      \item Menghitung perpindahan dari kecepatan dan selang waktu.
      \item Membuktikan kecepatan tetap pada 30 dan 60 bingkai per detik.
      \item Menyelaraskan posisi pecahan dengan \texttt{Rect} bilangan bulat.
    \end{itemize}
  \end{column}
\end{columns}
\end{frame}

% --------------------------------------------------------- peta
\petapertemuan{1}
  {Materi 2 \textperiodcentered\ Game Loop, Event Handling, dan Clock}
  {Materi 3 \textperiodcentered\ Gerak Objek, Vector2, Input, dan Delta Time}
  {Materi 4 \textperiodcentered\ Fisika Sederhana: Gravitasi, Lompatan, dan
   Pantulan}
  {Objek bergerak dengan kecepatan tetap dalam piksel per detik pada 30 dan
   60 FPS, posisi pecahan selaras dengan \texttt{Rect}, dan tabel pengukuran
   kecepatan terisi.}

% ========================================================= bagian 1
\section{Posisi dan Kecepatan sebagai Vektor}
\sectionframe{01}{Posisi dan Kecepatan sebagai Vektor}
  {Dua angka yang disatukan menjadi satu nilai}

\begin{frame}{Mengapa Pasangan Angka Perlu Disatukan}
\awalan{Posisi pada bidang datar dinyatakan oleh dua angka. Menyimpannya pada
dua variabel terpisah dapat dipakai, tetapi cepat menjadi merepotkan.}
\begin{itemize}
  \item Setiap operasi harus ditulis dua kali, dan setiap fungsi yang menerima
        posisi harus menerima dua argumen.
  \item Ketika jumlah objek bertambah, jumlah variabel bertambah dua kali
        lipat.
  \item Kekeliruan menukar pasangan angka---misalnya menambahkan kecepatan
        mendatar ke posisi tegak---sulit terlihat karena keduanya bertipe sama.
  \item \emph{Vector2} (\texttt{pygame.Vector2}) menyatukan kedua angka menjadi
        satu nilai; seluruh perhitungan dilakukan atas vektor, bukan atas
        komponennya.
  \item Penjumlahan dan perkalian dengan satu angka bekerja per komponen:
        $(200, 130) + (10, -5) = (210, 125)$ dan $(10, 5) \times 3 = (30, 15)$.
\end{itemize}
\vspace{0.16cm}
\begin{kotak}
Pada koordinat layar, sumbu tegak menuju ke bawah. Gerak ke atas karena itu
dinyatakan sebagai pengurangan nilai tegak. Tombol atas ditekan tetapi objek
bergerak ke bawah adalah kekeliruan tanda yang paling sering terjadi.
\end{kotak}
\end{frame}

\begin{frame}{Posisi, Kecepatan, dan Perpindahan sebagai Vektor}
\centering
\adjustbox{max width=0.82\linewidth,max height=0.55\textheight}{%
  % chapters/05/figures/vektor-kecepatan.tikz
% Posisi sebagai vektor, kecepatan sebagai vektor, dan perpindahan yang
% dihasilkan selama satu selang waktu. Perpindahan digambar segaris dengan
% kecepatan karena nilainya memang kecepatan dikalikan selang waktu.
\begin{tikzpicture}[
  panahposisi/.style={-{Stealth[length=2.4mm]}, very thick, draw=ubPrimary},
  panahkecepatan/.style={-{Stealth[length=2.4mm]}, very thick, draw=ubAccent},
  panahpindah/.style={-{Stealth[length=2.2mm]}, thick, draw=ubGreen, dashed},
  kisi/.style={draw=ubLine!50, very thin},
  label/.style={font=\scriptsize, align=left}
]

% Kisi latar sebagai gambaran koordinat layar.
\foreach \x in {0,1,...,10} \draw[kisi] (\x,0) -- (\x,-4);
\foreach \y in {0,-1,...,-4} \draw[kisi] (0,\y) -- (10,\y);

% Sumbu koordinat layar: x ke kanan, y ke bawah.
\draw[-{Stealth[length=2mm]}, thick, draw=ubLine] (0,0) -- (10.4,0)
  node[anchor=north west, font=\scriptsize, ubLine] {$x$};
\draw[-{Stealth[length=2mm]}, thick, draw=ubLine] (0,0) -- (0,-4.3)
  node[anchor=north east, font=\scriptsize, ubLine] {$y$};

% Titik-titik: posisi awal p, posisi sesudah satu selang waktu p', dan ujung
% vektor kecepatan. Perpindahan p ke p' segaris dengan kecepatan p ke v.
\coordinate (p) at (2,-2.6);
\coordinate (pp) at (3.36,-2.04);
\coordinate (v) at (5.4,-1.2);

\draw[panahposisi] (0,0) -- (p);
\fill[ubPrimary] (p) circle (1.8pt);
\node[label, ubPrimary, anchor=south east] at (1.9,-2.75)
  {posisi $p = (200;\ 130)$};

\draw[panahpindah] (p) -- (pp);
\fill[ubGreen] (pp) circle (1.6pt);

\draw[panahkecepatan] (p) -- (v);
\node[label, ubAccent, anchor=south west] at (4.2,-1.45)
  {kecepatan $v$ dalam piksel per detik};

\node[label, ubGreen, anchor=north west] at (2.35,-2.95)
  {perpindahan $v \cdot \Delta t$};

\node[label, ubGreen, anchor=west] at (3.5,-2.0) {$p'$};
\draw[ubGreen, very thin] (pp) -- (3.5,-2.0);
\end{tikzpicture}
}
\vspace{0.08cm}
\begin{itemize}
  \item Posisi $p = (200;\ 130)$ dan kecepatan $v$ masing-masing adalah satu
        nilai vektor, bukan dua angka yang kebetulan berdampingan.
  \item Perpindahan selama satu selang waktu adalah $v \cdot \Delta t$; arahnya
        segaris dengan kecepatan, panjangnya sebanding dengan selang waktu.
\end{itemize}
\sumber{Posisi, kecepatan, dan perpindahan sebagai vektor (Bab V buku ajar).}
\end{frame}

\begin{frame}{Operasi Vector2 yang Sering Dipakai}
\awalan{Operasi berikut disediakan Pygame dan dipakai sepanjang Blok I serta
Blok II.}
{\scriptsize
\renewcommand{\arraystretch}{1.10}
\begin{tabularx}{\textwidth}{@{}>{\raggedright\arraybackslash}p{4.7cm} X@{}}
\toprule
\textbf{Penulisan} & \textbf{Arti dan contoh} \\
\midrule
\texttt{pygame.Vector2(200, 130)} &
  Membuat vektor baru: 200 satuan mendatar dan 130 satuan tegak dari titik
  asal. \\
\texttt{position.x}, \texttt{position.y} &
  Membaca atau mengubah satu komponen; berguna untuk membatasi posisi pada tepi
  jendela. \\
\texttt{position + velocity * dt} &
  Menjumlahkan vektor dan mengalikannya dengan satu angka; perhitungan
  dilakukan per komponen secara otomatis. \\
\texttt{velocity.length()} &
  Panjang vektor, yaitu besarnya kecepatan dalam piksel per detik. \\
\texttt{velocity.length\_squared()} &
  Kuadrat panjang vektor; dipakai memeriksa vektor bernilai nol tanpa
  menghitung akar. \\
\texttt{direction.normalize()} &
  Vektor baru yang arahnya sama dengan panjang satu; vektor asalnya tidak
  berubah. \\
\texttt{a.distance\_to(b)} &
  Jarak antara dua titik; berguna untuk mengukur jarak pemain terhadap objek
  lain. \\
\bottomrule
\end{tabularx}}
\vspace{0.10cm}
\sumber{\texttt{normalize()} mengembalikan vektor baru; \texttt{normalize\_ip()}
mengubah vektor asalnya. Menulis \texttt{direction.normalize()} tanpa
menyimpan hasilnya tidak berpengaruh apa pun.}
\end{frame}

% ========================================================= bagian 2
\section{Membaca Masukan yang Berlangsung Terus}
\sectionframe{02}{Membaca Masukan yang Berlangsung Terus}
  {Keadaan tombol, bukan peristiwa sekali}

\begin{frame}{Keadaan Seluruh Tombol}
\awalan{\texttt{pygame.key.get\_pressed()} mengembalikan keadaan seluruh tombol
pada saat itu, bukan peristiwa.}
\begin{itemize}
  \item Pemeriksaan dilakukan tanpa menunggu peristiwa, sehingga hasilnya selalu
        menggambarkan keadaan terkini.
  \item Bila tombol masih ditahan pada bingkai ini, nilainya benar; bila sudah
        dilepas, nilainya salah. Karena itu cara ini dipakai untuk gerak yang
        berlangsung terus.
  \item Dua tombol dihubungkan dengan \texttt{or}, sehingga pemain dapat memakai
        tombol panah maupun huruf WASD.
  \item Peristiwa \texttt{KEYDOWN} hanya terjadi sekali saat tombol ditekan dan
        tidak memadai untuk gerak selama tombol ditahan.
  \item Vektor arah dibaca sekali setiap bingkai, pada awal tahap pembaruan,
        lalu dipakai menghitung kecepatan.
\end{itemize}
\vspace{0.16cm}
\begin{kotak}
Gerak yang dihasilkan disebut gerak langsung: kecepatan penuh begitu tombol
ditekan, dan langsung berhenti begitu dilepas. Percepatan dan perlambatan
ditambahkan pada Bab VI.
\end{kotak}
\end{frame}

\begin{frame}{Menyusun Vektor Arah dan Menormalkannya}
\centering
\adjustbox{max width=0.86\linewidth,max height=0.37\textheight}{%
  % chapters/05/figures/normalisasi-diagonal.tikz
% Gerak diagonal tanpa dan dengan penormalan vektor arah.
\begin{tikzpicture}[
  panah/.style={-{Stealth[length=2.2mm]}, very thick},
  kisi/.style={draw=ubLine!60, very thin},
  judul/.style={font=\footnotesize\bfseries, ubPrimary},
  label/.style={font=\scriptsize, align=center}
]

% ---------- Tanpa penormalan ----------
\begin{scope}
  \foreach \x in {0,1,2} \draw[kisi] (\x,0) -- (\x,-2);
  \foreach \y in {0,-1,-2} \draw[kisi] (0,\y) -- (2,\y);
  \draw[panah, ubRed] (0,0) -- (1.7,-1.7);
  \node[label, ubRed, anchor=north] at (1.0,-2.15)
    {panjang $\sqrt{2} \approx 1{,}414$};
  \node[judul, anchor=south] at (1.0,0.25) {tanpa penormalan};
  \node[label, anchor=north, text width=42mm] at (1.0,-2.75)
    {kecepatan menjadi $283$ piksel/detik, yaitu $41$ persen lebih cepat};
\end{scope}

% ---------- Dengan penormalan ----------
\begin{scope}[xshift=7.2cm]
  \foreach \x in {0,1,2} \draw[kisi] (\x,0) -- (\x,-2);
  \foreach \y in {0,-1,-2} \draw[kisi] (0,\y) -- (2,\y);
  \draw[panah, ubGreen] (0,0) -- (1.2,-1.2);
  \node[label, ubGreen, anchor=north] at (1.0,-2.15)
    {panjang $1$};
  \node[judul, anchor=south] at (1.0,0.25) {dengan penormalan};
  \node[label, anchor=north, text width=42mm] at (1.0,-2.75)
    {kecepatan tetap $200$ piksel/detik, sama dengan gerak lurus};
\end{scope}
\end{tikzpicture}
}
\vspace{0.04cm}
\begin{itemize}
  \item Dua tombol bersamaan, misalnya kanan dan bawah, menghasilkan vektor
        $(1, 1)$ dengan panjang $\sqrt{1^2 + 1^2} = \sqrt{2} \approx 1{,}414$.
  \item Akibatnya objek bergerak 41 persen lebih cepat ketika menempuh arah
        diagonal: $200 \times 1{,}414 = 282{,}8$ piksel per detik.
  \item Penanganannya disebut \emph{normalisasi}: panjang vektor dipaksa menjadi
        satu sebelum dikalikan kecepatan, yaitu
        \texttt{direction.normalize() * speed}.
  \item Pemeriksaan \texttt{length\_squared() > 0} dilakukan lebih dahulu karena
        menormalkan vektor bernilai nol tidak bermakna.
\end{itemize}
\sumber{Akibat gerak diagonal tanpa dan dengan penormalan vektor arah (Bab V).}
\end{frame}

\begin{frame}{Urutan Pembacaan yang Benar}
\begin{itemize}
  \item Vektor arah dibaca \textbf{sekali} pada setiap bingkai, yaitu pada awal
        tahap pembaruan, lalu dipakai untuk menghitung kecepatan.
  \item Membaca tombol berulang kali di dalam satu bingkai tidak salah, tetapi
        tidak ada gunanya: keadaan tombol tidak berubah di tengah bingkai.
  \item Yang perlu dihindari adalah membaca tombol pada tahap penayangan,
        misalnya di dalam fungsi penggambaran. Keadaan yang digambar tidak
        selalu sama dengan keadaan yang dihitung.
  \item Kebiasaan yang benar: masukan dibaca pada tahap masukan, disimpan pada
        variabel, lalu dipakai oleh tahap-tahap berikutnya.
\end{itemize}
\vspace{0.16cm}
\begin{kiat}
Untuk memeriksa apakah vektor bernilai nol, pakailah
\texttt{length\_squared()}, bukan \texttt{length()}. Keduanya memberi jawaban
yang sama untuk keperluan itu, tetapi kuadrat panjang lebih murah dihitung
karena tidak menuntut operasi akar.
\end{kiat}
\end{frame}

% ========================================================= bagian 3
\section{Gerak Berbasis Selang Waktu}
\sectionframe{03}{Gerak Berbasis Selang Waktu}
  {Menghubungkan kecepatan dengan waktu yang terukur}

\begin{frame}{Perpindahan Selama Satu Bingkai}
\awalan{Perpindahan dihitung dari kecepatan dan selang waktu yang benar-benar
berlangsung, yaitu nilai yang dikembalikan \texttt{clock.tick()}.}
\[
\Delta t = \frac{\text{clock.tick(FPS)}}{1000}\ \text{detik},
\qquad
\vec{p}_{\text{baru}} = \vec{p} + \vec{v}\,\Delta t .
\]
\begin{itemize}
  \item Kecepatan dinyatakan dalam piksel per detik dan selang waktu dalam
        detik, sehingga hasil kalinya adalah piksel.
  \item Apabila selang waktunya dua kali lebih panjang, perpindahannya dua kali
        lebih jauh, dan jarak yang ditempuh dalam satu detik tetap sama.
  \item Urutan perhitungan mengikuti kaidah Euler: kecepatan dihitung lebih
        dahulu, baru posisi diperbarui.
  \item Baris \texttt{position = position + velocity * delta\_seconds} inilah
        yang membuat kecepatan tidak bergantung pada laju bingkai.
\end{itemize}
\vspace{0.14cm}
\begin{kotak}
Selang waktu bingkai pertama bernilai nol, sehingga bingkai itu belum
menghasilkan perpindahan. Nilai itu tetap harus dipakai apa adanya, bukan
diganti dengan jatah waktu yang diharapkan.
\end{kotak}
\end{frame}

\begin{frame}{Membuktikan Kecepatan Tetap melalui Pengukuran}
\awalan{Setiap bingkai program mencatat perubahan posisi yang sesungguhnya,
lalu membagi jumlah jarak dengan lama objek benar-benar bergerak.}
{\scriptsize
\renewcommand{\arraystretch}{1.22}
\begin{tabularx}{\textwidth}{@{}>{\centering\arraybackslash}p{1.9cm}
    >{\centering\arraybackslash}p{1.5cm}
    >{\centering\arraybackslash}p{2.7cm}
    >{\centering\arraybackslash}p{2.5cm} X@{}}
\toprule
\textbf{Batas laju} & \textbf{Bingkai} & \textbf{Waktu berjalan} &
\textbf{Jarak tempuh} & \textbf{Kecepatan terukur} \\
\midrule
30 FPS  & 18 & 0,61 detik & 115,4 piksel & 200,0 piksel/detik \\
60 FPS  & 36 & 0,60 detik & 116,8 piksel & 200,0 piksel/detik \\
120 FPS & 72 & 0,62 detik & 122,4 piksel & 200,0 piksel/detik \\
\bottomrule
\end{tabularx}}
\vspace{0.16cm}
\begin{itemize}
  \item Kolom terakhir sama pada ketiga laju: pada 30 FPS objek berpindah lebih
        jauh setiap bingkainya, pada 120 FPS lebih pendek, tetapi jarak per
        detiknya sama.
  \item Jarak tempuh boleh berbeda sedikit karena \texttt{clock.tick()} hanya
        menjamin laju tidak \emph{melampaui} batas, bukan bahwa setiap bingkai
        tepat sepanjang jatahnya.
  \item Yang harus sama adalah \textbf{kecepatan}, bukan jarak tempuh selama
        selang waktu yang diukur secara kasar.
\end{itemize}
\sumber{Pengukuran penulis pada kecepatan 200 piksel per detik, arah lurus,
tanpa jendela; angka pada komputer Anda akan berbeda.}
\end{frame}

\begin{frame}{Mengapa Cara Ini Perlu Diuji Sendiri}
\awalan{Anggapan ``jelas kecepatannya sama'' berbahaya, karena tiga kekeliruan
berikut tidak menimbulkan pesan kesalahan apa pun.}
\begin{itemize}
  \item \texttt{delta\_seconds} lupa dipakai sebagai pengali: program tetap
        berjalan, yang terlihat hanyalah gerak yang berbeda-beda menurut
        komputer.
  \item \texttt{delta\_seconds} dikalikan dua kali, misalnya karena kecepatan
        sudah dinyatakan per bingkai: objek tetap bergerak dan hanya terasa
        terlalu cepat.
  \item Penormalan vektor arah terlupa: gerak diagonal tetap tampak wajar sampai
        pemain menyadari bahwa berlari diagonal lebih cepat daripada berlari
        lurus.
\end{itemize}
\vspace{0.16cm}
\begin{kiat}
Biasakan menjalankan program gerak pada dua laju bingkai yang berbeda, misalnya
30 dan 60, setiap kali selesai menulisnya. Apabila kecepatan objek terasa
berbeda, ada yang keliru pada perhitungan perpindahannya.
\end{kiat}
\end{frame}

% ========================================================= bagian 4
\section{Menyinkronkan Posisi Pecahan dengan Rect}
\sectionframe{04}{Menyinkronkan Posisi Pecahan\\ dengan Rect}
  {Dua tipe data, dua tugas yang berbeda}

\begin{frame}[fragile]{Dua Tipe Data yang Berbeda}
\awalan{Posisi gerak berjalan lancar sebagai pecahan, tetapi persoalan muncul
ketika posisi itu hendak dipakai untuk menggambar.}
\begin{itemize}
  \item \texttt{pygame.Vector2} menyimpan nilai pecahan, misalnya 312,47.
  \item \emph{Rect} menyimpan bilangan bulat: piksel ke-312 tepat berada di
        tempatnya, dan tidak ada piksel ke-312,47.
  \item Posisi pecahan karena itu diselaraskan setiap bingkai, tepat sebelum
        penggambaran.
\end{itemize}
\vspace{0.10cm}
\lstinputlisting[language=Python,firstline=197,lastline=199]
  {../buku_ajar/code/python/ch05_gerak.py}
\kodejelas{Pembulatan, bukan pemotongan, agar gerak tidak tampak tersendat.
\texttt{player\_rect} hanyalah salinan yang dibulatkan untuk menggambar dan
untuk memeriksa tabrakan pada Bab VII; posisi sesungguhnya tetap disimpan pada
\texttt{position}.}
\end{frame}

\begin{frame}{Posisi Tetap Disimpan sebagai Pecahan}
\awalan{Aturan yang perlu dipegang adalah pembagian tugas antara posisi yang
sebenarnya dan posisi yang digambar.}
\begin{itemize}
  \item \textbf{Posisi sebenarnya disimpan sebagai pecahan} (\texttt{Vector2}),
        karena hanya dengan begitu gerak dapat berlangsung halus. Pada kecepatan
        20 piksel per detik dan 60 bingkai per detik, perpindahan setiap
        bingkainya hanya $20 \div 60 = 0{,}333$ piksel; nilai itu harus
        tersimpan agar dapat terakumulasi.
  \item \textbf{\texttt{Rect} dipakai untuk menggambar dan memeriksa tabrakan},
        dan diselaraskan dari posisi pecahan setiap bingkai. \texttt{Rect} bukan
        tempat menyimpan posisi.
  \item Bila aturan itu dilanggar, objek yang posisinya bilangan bulat tidak
        bergerak sama sekali pada kecepatan rendah: perpindahan 0,333 piksel
        terpotong menjadi nol pada setiap bingkai, dan nilai yang terpotong itu
        hilang untuk selamanya.
  \item Gejala itu tidak muncul pada kecepatan tinggi. Pada 200 piksel per detik
        perpindahannya 3,33 piksel, sehingga kekeliruannya tidak terlihat sampai
        kecepatan diturunkan untuk permainan yang sesungguhnya.
\end{itemize}
\end{frame}

\begin{frame}{Membulatkan atau Memotong}
\awalan{Pilihan antara \texttt{round()} dan \texttt{int()} memengaruhi letak
objek yang terlihat.}
\begin{itemize}
  \item \texttt{int(312.9)} menghasilkan 312 (selalu ke arah nol), sedangkan
        \texttt{round(312.9)} menghasilkan 313 (ke bilangan bulat terdekat).
  \item Karena posisi yang benar adalah 312,9, nilai yang lebih dekat adalah
        313. Memakai \texttt{int()} membuat objek tampak tertinggal hampir satu
        piksel secara sistematis.
  \item Selisih itu tidak terlihat pada objek tunggal, tetapi menjadi gerak yang
        tersendat ketika beberapa objek bergerak bersama-sama.
\end{itemize}
\vspace{0.14cm}
\begin{kotak}
Memisahkan ``posisi yang sebenarnya'' dari ``posisi yang digambar'' membawa
tiga manfaat: gerak tetap halus pada kecepatan rendah, \texttt{Rect} untuk
tabrakan tersedia setiap bingkai tanpa mengorbankan ketelitian, dan perubahan
laju bingkai tidak mengubah posisi akhir.
\end{kotak}
\end{frame}

% ========================================================= bagian 5
\section{Program Gerak Lengkap}
\sectionframe{05}{Program Gerak Lengkap}
  {Masukan, pembaruan, penyelarasan, dan penayangan}

\begin{frame}[fragile]{Tahap Masukan: read\_direction}
\awalan{Fungsi ini mengembalikan vektor arah berdasarkan tombol yang sedang
ditahan.}
\lstinputlisting[language=Python,firstline=80,lastline=90]
  {../buku_ajar/code/python/ch05_gerak.py}
\kodejelas{Vektor yang dikembalikan belum dinormalkan: panjangnya satu bila satu
tombol ditekan, dan $\sqrt{2}$ bila dua tombol ditekan bersamaan. Penormalan
dilakukan di tempat lain agar fungsi ini tetap sederhana dan hanya menangani
masukan.}
\end{frame}

\begin{frame}[fragile]{Tahap Pembaruan: update\_position}
\lstinputlisting[language=Python,firstline=100,lastline=108]
  {../buku_ajar/code/python/ch05_gerak.py}
\kodejelas{Arah dinormalkan lalu dikalikan kecepatan; posisi diperbarui dengan
perpindahan selama satu selang waktu. Ketiga pekerjaan---penormalan,
pembaruan posisi, dan pembatasan pada tepi jendela---berkaitan erat sehingga
diletakkan pada satu tempat.}
\vspace{0.10cm}
\begin{kotak}
Pembatasan tepi jendela dilakukan dengan memotong nilai posisi, bukan dengan
memeriksa tabrakan. Pemeriksaan tabrakan yang sesungguhnya baru dibahas pada
Bab VII.
\end{kotak}
\end{frame}

\begin{frame}{Percobaan yang Disarankan}
\awalan{Setiap perubahan kecil berikut dapat diamati langsung pada jendela
program.}
{\scriptsize
\renewcommand{\arraystretch}{1.08}
\begin{tabularx}{\textwidth}{@{}>{\raggedright\arraybackslash}p{5.4cm} X@{}}
\toprule
\textbf{Yang diubah dan caranya} & \textbf{Yang sebaiknya Anda amati} \\
\midrule
Laju bingkai melalui \texttt{--fps}. &
  Kecepatan terukur tetap sama, sedangkan perpindahan setiap bingkai berubah.
  Bandingkan nilai \texttt{delta\_seconds}. \\
Kecepatan melalui \texttt{--speed}. &
  Kecepatan terukur berubah sebanding. Pada kecepatan sangat rendah, misalnya
  5 piksel per detik, gerak menjadi tersendat karena pembulatan. \\
Arah diagonal dengan menekan dua tombol arah sekaligus. &
  Kecepatan terukur tetap sama dengan gerak lurus; inilah bukti penormalan
  bekerja. \\
Baris penormalan dengan menghapus \texttt{normalize()}. &
  Kecepatan diagonal naik menjadi sekitar 283 piksel per detik. Kembalikan baris
  itu sesudah percobaan. \\
Pembulatan dengan mengganti \texttt{round()} menjadi \texttt{int()}. &
  Selisih letak hampir satu piksel; baris \texttt{Rect} pada jendela
  memperlihatkan nilainya berbeda. \\
\bottomrule
\end{tabularx}}
\vspace{0.12cm}
\sumber{Kembalikan program ke keadaan semula sesudah bereksperimen, dan simpan
keadaan percobaan pada repositori Git.}
\end{frame}

% ========================================================= bagian 6
\section{Yang Belum Ada pada Gerak Ini}
\sectionframe{06}{Yang Belum Ada pada Gerak Ini}
  {Batas kemampuan program dan apa yang ditambahkan}

\begin{frame}{Kecepatan yang Selalu Tetap}
\awalan{Kecepatan pada bab ini hanya berubah karena tombol ditekan atau
dilepas; selama ditahan nilainya tetap. Gerak semacam itu disebut gerak lurus
beraturan.}
\begin{itemize}
  \item Sebagian besar permainan tidak cukup dengan gerak itu: benda yang jatuh
        makin lama makin cepat, bola yang memantul kehilangan sebagian
        kecepatannya, dan karakter yang melompat melambat mendekati titik
        tertinggi.
  \item Ketiga gejala itu berasal dari \textbf{percepatan}, yaitu perubahan
        kecepatan terhadap waktu.
  \item Percepatan ditambahkan pada tahap pembaruan, tepat sebelum kecepatan
        dipakai untuk memperbarui posisi.
\end{itemize}
\[
\vec{v} \leftarrow \vec{v} + \vec{a}\,\Delta t,
\qquad
\vec{p} \leftarrow \vec{p} + \vec{v}\,\Delta t .
\]
\begin{kotak}
Baris pertama itulah yang ditambahkan pada Bab VI, beserta cara memilih nilai
percepatan agar gravitasi terasa wajar dan lompatan dapat dikendalikan.
\end{kotak}
\end{frame}

\begin{frame}{Objek yang Belum Dapat Bertabrakan dan Belum Menjadi Sprite}
\begin{itemize}
  \item Objek pada bab ini bergerak menembus apa pun. Ia hanya dibatasi tepi
        jendela, dan pembatasan itu pun dilakukan dengan memotong nilai posisi,
        bukan dengan memeriksa tabrakan.
  \item Pada Bab VII, posisi pecahan dan \texttt{Rect} yang telah diselaraskan
        dipakai untuk memeriksa kapan dua objek bersentuhan. Pemeriksaan
        tabrakan memerlukan \texttt{Rect} bilangan bulat, sedangkan ketelitian
        gerak memerlukan posisi pecahan.
  \item Program ini menggerakkan satu objek yang disimpan pada beberapa
        variabel: \texttt{position}, \texttt{velocity}, dan \texttt{player\_rect}.
        Cara itu memadai untuk satu objek, tetapi tidak terkelola untuk sepuluh
        atau seratus objek.
  \item Kelas \emph{Sprite} menyatukan gambar, posisi, dan perilaku sebuah objek;
        kelas \emph{Group} mengelola sekumpulan sprite sekaligus. Keduanya
        dibahas pada Bab VII.
  \item Sesudah itu, menambahkan objek baru cukup dengan membuat satu objek
        lagi, bukan dengan menambah tiga variabel baru untuk setiap objek.
\end{itemize}
\end{frame}

\begin{frame}{Ringkasan Perkembangan}
\awalan{Dua hal dari bab ini akan dipakai terus sampai akhir buku: posisi
sebagai pecahan yang diperbarui setiap bingkai, dan perpindahan sebagai hasil
kali kecepatan dengan selang waktu.}
{\scriptsize
\renewcommand{\arraystretch}{1.14}
\begin{tabularx}{\textwidth}{@{}>{\centering\arraybackslash}p{1.7cm}
    >{\raggedright\arraybackslash}p{5.0cm} X@{}}
\toprule
\textbf{Bab} & \textbf{Yang ditambahkan} & \textbf{Akibatnya} \\
\midrule
Bab VI & Percepatan: gravitasi, impuls lompatan, pantulan &
  Kecepatan tidak lagi tetap, melainkan berubah setiap bingkai menurut
  percepatan yang bekerja. \\
Bab VII & \emph{Sprite}, \emph{Group}, dan deteksi tabrakan &
  Objek menjadi satu kesatuan beserta gambarnya, dan perpindahannya dapat
  terhenti karena bertabrakan. \\
Bab VIII--XI & Pengendalian NPC &
  Kecepatan tidak lagi berasal dari tombol, melainkan dari keputusan yang
  diambil program sendiri. \\
Bab XIII & Fisika mesin pada Unity &
  Perhitungan percepatan dan tabrakan diserahkan kepada mesin, sedangkan kode
  Anda mengatur pengaturan dan menanggapi peristiwanya. \\
\bottomrule
\end{tabularx}}
\vspace{0.14cm}
\sumber{Bagian yang ditambahkan pada gerak objek pada bab-bab berikutnya
(Bab V).}
\end{frame}

% --------------------------------------------------------- aktivitas
\begin{frame}{Aktivitas di Kelas \textperiodcentered\ Praktikum 03}
\setbeamerfont{itemize/enumerate body}{size=\footnotesize}
\awalan{Pengalaman belajar minggu 3 pada RPS, dikerjakan berurutan di kelas lalu
diselesaikan di rumah.}
\begin{enumerate}
  \item Ketik ulang \texttt{ch05\_gerak.py} --- jangan salin-tempel --- lalu
        gerakkan persegi dengan tombol WASD dan dengan tombol panah.
  \item Amati dua baris posisi pada jendela: nilai pecahan berubah halus,
        sedangkan nilai \texttt{Rect} melompat satu piksel setiap kali.
  \item Ukur kecepatan pada 30 FPS:
        \texttt{python ch05\_gerak.py --fps 30 --frames 18 --auto-direction
        kanan}, lalu catat kecepatan terukurnya.
  \item Ulangi pada 60 dan 120 FPS dengan \texttt{--frames 36} dan
        \texttt{--frames 72}; bandingkan kecepatan terukurnya.
  \item Uji \texttt{--auto-direction diagonal}; hapus pemanggilan
        \texttt{normalize()}, catat perbedaannya, lalu kembalikan baris itu.
  \item Tambahkan objek kedua berkecepatan 120 piksel per detik dengan tombol
        IJKL, tulis catatan praktikum, dan simpan pekerjaan ke Git.
\end{enumerate}
\vspace{0.10cm}
{\footnotesize\color{ubAccent}\itshape Rincian lengkap dan lembar kerja:
Praktikum 03 pada Bab V buku ajar. Kumpulkan melalui repositori Git pribadi.}
\end{frame}

% --------------------------------------------------------- rangkuman
\begin{frame}{Rangkuman}
\begin{itemize}
  \item \texttt{pygame.Vector2} menyatukan dua komponen posisi atau kecepatan
        menjadi satu nilai; penjumlahan dan perkalian dengan satu angka bekerja
        per komponen.
  \item Pada koordinat layar sumbu tegak menuju ke bawah, sehingga gerak ke atas
        dinyatakan sebagai pengurangan nilai tegak.
  \item \texttt{pygame.key.get\_pressed()} menghasilkan gerak yang berlangsung
        selama tombol ditahan; peristiwa \texttt{KEYDOWN} hanya terjadi sekali.
  \item Vektor arah dinormalkan menjadi panjang satu sebelum dikalikan kecepatan;
        tanpa itu gerak diagonal menjadi 41 persen lebih cepat.
  \item $\Delta t = \text{clock.tick(FPS)}/1000$ dan
        $p_{\text{baru}} = p + v\,\Delta t$; kecepatan terukur tetap 200,0
        piksel per detik pada 30, 60, dan 120 FPS.
  \item Posisi disimpan sebagai pecahan, sedangkan \texttt{Rect} dipakai untuk
        menggambar dan memeriksa tabrakan; \texttt{round()} dipakai, bukan
        \texttt{int()}.
\end{itemize}
\end{frame}

% --------------------------------------------------------- cek
\begin{frame}{Cek Pemahaman}
\begin{enumerate}
  \item Mengapa vektor arah $(1, 1)$ perlu dinormalkan sebelum dikalikan
        kecepatan? Sertakan perhitungan singkat untuk mendukung jawaban Anda.
  \item Mengapa \texttt{delta\_seconds} harus menjadi pengali perpindahan, dan
        gejala apa yang muncul apabila pengali itu terlupa?
  \item Mengapa posisi objek disimpan sebagai bilangan pecahan, sedangkan
        \texttt{Rect} bertipe bilangan bulat? Apa akibatnya pada kecepatan
        20 piksel per detik?
\end{enumerate}
\vspace{0.35cm}
{\footnotesize\color{ubMuted}\itshape Jawab tanpa membuka catatan, lalu
bandingkan dengan pasangan di sebelah Anda.}
\end{frame}

% --------------------------------------------------------- penutup
\penutup{Tugas Praktikum 03}
  {Pada Materi 4 kecepatan tidak lagi tetap: percepatan \mbox{gravitasi},
   impuls lompatan, dan pantulan dinding dihitung dengan integrasi Euler yang
   sama, lalu \mbox{diuji} pada nilai \mbox{gravitasi} dan koefisien pantulan
   yang berbeda.}

\end{document}
